\scanpage{187}
\begin{center}- 63 -\end{center}
\hfill IX
\uline{Proposition} 5.7. \uline{Supposons} $A^{\unsure{o|0}}$ \uline{$\ell$-divisible}, $S$ \uline{hensélien}. \uline{Soit} $F_K$ \uline{un}
\uline{faisceau $\ell$-adique constructible sur} $K$ \uline{qui soit} ``\uline{non ramifié}'' \uline{i.e. qui}
\uline{se prolonge en un faisceau $\ell$-adique constant tordu} $F$ \uline{sur} $S$, \uline{et soit}
$u:F_K\longrightarrow T_\ell(A_K)$ \uline{un homomorphisme de pro-groupes}, \uline{alors} $u$ \uline{se factorise}
\uline{en un homomorphisme de} $F$ \uline{dans} $T_\ell(A_K)^{\mathrm f}$.

En effet, on aura
\[
(5.7.1)\qquad u=(u(\nu))_{\nu\geq0},\quad\text{avec }u(\nu):F_K(\nu)\longrightarrow T_\ell(A_K)(\nu)={}_{\ell^\nu}A_K\quad .
\]
Par la propriété universelle de $A$, comme $F(\nu)$ est étale donc lisse sur
$S$, le morphisme $u(\nu)$ provient d'un homomorphisme
\[
v(\nu):F(\nu)\longrightarrow{}_{\ell^\nu}A\quad ,
\]
et par définition de la partie fixe, cet homomorphisme se factorise en
\[
v(\nu):F(\nu)\longrightarrow({}_{\ell^\nu}A)^{\mathrm f}\quad .
\]
Comme l'image du morphisme induit sur les fibres spéciales tombe dans
la partie infiniment $\ell$-divisible de $A_{\unsure{o|0}}$, on conclut que $v(\nu)$ se factorise
même par $({}_{\ell^\nu}A^{\unsure{o|0}})^{\mathrm f}$. Par passage à la limite sur $\nu$, on trouve un homomorphisme
\[
v:F\longrightarrow T_\ell(A^{\unsure{o|0}})^{\mathrm f}\quad ,
\]
et le morphisme induit sur les fibres génériques donne évidemment la
factorisation cherchée.

\uline{Remarques} 5.7.2. Lorsque $\ell\ne p$, il est clair que la propriété 5.7 du
sous-pro-groupe $T_\ell(A_K)^{\mathrm f}$ de $T_\ell(A_K)$ caractérise celui-ci de façon unique,
comme le plus grand sous-faisceau $\ell$-adique de $T_\ell(A_K)$ qui soit ``non ramifié''.
Lorsque $\ell=p$, l'énoncé 5.7 cesse d'être bien raisonnable, et on va donner

\margin{381}
\margin{381}
\margin{186}
